\ifdefined\gkver\else\def\gkver{student}\fi
\documentclass[\gkver]{gaokaozhenti}
\begin{document}

\chapter{2022年全国乙卷}
%% paperType: 理综物理部分

\begin{enumerate}
\item
%% number: 1
%% paperName: 2022年全国高考乙卷第 1 题 6 分
%% typeId: 1
%% type: 单选题
%% chapter: 曲线运动
%% point: 人造地球卫星
%% method: 无
%% score: 6
%% degree: 750
%% duplicateId: 0
%% body:
2022 年 3 月，中国航天员翟志刚、王亚平、叶光富在离地球表面约 $400\Ukm$ 的“天宫二号”空间站上通过天地连线，为同学们上了一堂精彩的科学课。通过直播画面可以看到，在近地圆轨道上飞行的“天宫二号”中，航天员可以自由地漂浮，这表明他们\xzanswer{C}

\fourchoices[answer=C]
{所受地球引力的大小近似为零}
{所受地球引力与飞船对其作用力两者的合力近似为零}
{所受地球引力的大小与其随飞船运动所需向心力的大小近似相等}
{在地球表面上所受引力的大小小于其随飞船运动所需向心力的大小}

%% answer: C
%% memo:
\memoanswer{
ABC．航天员在空间站中所受万有引力完全提供做圆周运动的向心力，飞船对其作用力等于零，故 C 正确，AB 错误；

D．根据万有引力公式 $F = G \frac{Mm}{r^{2}}$ 可知，在地球表面上所受引力的大小大于在飞船所受的万有引力大小，因此地球表面引力大于其随飞船运动所需向心力的大小，故 D 错误。

故选 C。
}

\item
%% number: 2
%% paperName: 2022年全国高考乙卷第 2 题 6 分
%% typeId: 1
%% type: 单选题
%% chapter: 运动和力
%% point: 牛顿第二定律
%% method: 无
%% score: 6
%% degree: 650
%% duplicateId: 0
%% body:
如图，一不可伸长轻绳两端各连接一质量为 $m$ 的小球，初始时整个系统静置于光滑水平桌面上，两球间的距离等于绳长 $L$。一大小为 $F$ 的水平恒力作用在轻绳的中点，方向与两球连线垂直。当两球运动至二者相距 $\frac{3}{5}L$ 时，它们加速度的大小均为\xzanswer{A}

\onepicture[w=3.0cm]{figs/02.png}

\fourchoices[answer=A]
{$\frac{5F}{8m}$}
{$\frac{2F}{5m}$}
{$\frac{3F}{8m}$}
{$\frac{3F}{10m}$}

%% answer: A
%% memo:
\memoanswer{
当两球运动至二者相距 $\frac{3}{5}L$ 时，如图所示由几何关系可知

\onepicture[w=5.5cm]{figs/02_an.png}

$\sin\theta=\frac{\frac{3L}{10}}{\frac{L}{2}}=\frac{3}{5}$

设绳子拉力为 $T$，水平方向有 $2T\cos\theta=F$

解得 $T=\frac{5}{8}F$

对任意小球由牛顿第二定律可得：$T=ma$

解得：$a=\frac{5F}{8m}$

故 A 正确，BCD 错误。

故选 A。
}

\item
%% number: 3
%% paperName: 2022年全国高考乙卷第 3 题 6 分
%% typeId: 1
%% type: 单选题
%% chapter: 功和能
%% point: 机械能守恒定律
%% method: 无
%% score: 6
%% degree: 450
%% duplicateId: 0
%% body:
固定于竖直平面内的光滑大圆环上套有一个小环，小环从大圆环顶端 P 点由静止开始自由下滑，在下滑过程中，小环的速率正比于\xzanswer{C}

\onepicture[w=4.0cm]{\includesvg{figs/03}}

\fourchoices[answer=C]
{它滑过的弧长}
{它下降的高度}
{它到 P 点的距离}
{它与 P 点的连线扫过的面积}

%% answer: C
%% memo:
\memoanswer{
如图所示

\onepicture[w=4.5cm]{figs/03_an.png}

设圆环下降的高度为 $h$，圆环的半径为 $R$，它到 P 点的距离为 $L$，根据机械能守恒定律得

$mgh=\frac{1}{2}mv^{2}$

由几何关系可得 $h = L\sin\theta$

$\sin\theta = \frac{L}{2R}$

联立可得 $h = \frac{L^{2}}{2R}$

可得 $v = L\sqrt{\frac{g}{R}}$

故 C 正确，ABD 错误。

故选 C。
}

\item
%% number: 4
%% paperName: 2022年全国高考乙卷第 4 题 6 分
%% typeId: 1
%% type: 单选题
%% chapter: 光学
%% point: 光子说
%% method: 建立模型
%% score: 6
%% degree: 450
%% duplicateId: 0
%% body:
一点光源以 $113\UW$ 的功率向周围所有方向均匀地辐射波长约为 $6\times10^{-7}\Um$ 的光，在离点光源距离为 $R$ 处每秒垂直通过每平方米的光子数为 $3\times10^{14}$ 个。普朗克常量为 $h = 6.63\times10^{-34}\UJcdots$。$R$ 约为\xzanswer{B}

\fourchoices[answer=B]
{$1\times10^{2}\Um$}
{$3\times10^{2}\Um$}
{$6\times10^{2}\Um$}
{$9\times10^{2}\Um$}

%% answer: B
%% memo:
\memoanswer{
一个光子的能量为 $E = h\nu$

$N$ 为光的频率，光的波长与频率有以下关系 $c = \lambda v$

光源每秒发出的光子的个数为 $n = \frac{P}{h\nu} = \frac{P\lambda}{hc}$

$P$ 为光源的功率，光子以球面波的形式传播，那么以光源为原点的球面上的光子数相同，此时距光源的距离为 $R$ 处，每秒垂直通过每平方米的光子数为 $3\times10^{14}$ 个，那么此处的球面的表面积为 $S = 4\pi R^{2}$

则 $\frac{n}{S} = 3\times10^{14}$

联立以上各式解得 $R \approx 3\times10^{2}\Um$

故选 B。
}

\item
%% number: 5
%% paperName: 2022年全国高考乙卷第 5 题 6 分
%% typeId: 2
%% type: 多选题
%% chapter: 磁场
%% point: 磁场的叠加
%% method: 无
%% score: 6
%% degree: 450
%% duplicateId: 0
%% body:
安装适当的软件后，利用智能手机中的磁传感器可以测量磁感应强度 $B$。如图，在手机上建立直角坐标系，手机显示屏所在平面为 $xOy$ 面。某同学在某地对地磁场进行了四次测量，每次测量时 $y$ 轴指向不同方向而 $z$ 轴正向保持竖直向上。根据表中测量结果可推知\xzanswer{BC}

\begin{table}[htbp]
\centering
\begin{tblr}{|c|c|c|c|}
\hline
测量序号 & $B_{x}/\UuT[a]$ & $B_{y}/\UuT[a]$ & $B_{z}/\UuT[a]$ \\
\hline
1 & 0 & 21 & $-$45 \\
\hline
2 & 0 & $-$20 & $-$46 \\
\hline
3 & 21 & 0 & $-$45 \\
\hline
4 & $-$21 & 0 & $-$45 \\
\hline
\end{tblr}
\end{table}

\onepicture[w=4.5cm]{figs/05.png}

\fourchoices[answer=BC]
{测量地点位于南半球}
{当地的地磁场大小约为 $50\UuT$}
{第 2 次测量时 $y$ 轴正向指向南方}
{第 3 次测量时 $y$ 轴正向指向东方}

%% answer: BC
%% memo:
\memoanswer{
A. 如图所示

\onepicture[w=4.5cm]{figs/05_an.png}

地球可视为一个磁偶极，磁南极大致指向地理北极附近，磁北极大致指向地理南极附近。通过这两个磁极的假想直线（磁轴）与地球的自转轴大约成 11.3 度的倾斜。由表中 $z$ 轴数据可看出 $z$ 轴的磁场竖直向下，则测量地点应位于北半球，A 错误；

B. 磁感应强度为矢量，故由表格可看出此处的磁感应强度大致为 $B = \sqrt{B_{x}^{2} + B_{z}^{2}} = \sqrt{B_{y}^{2} + B_{z}^{2}}$

计算得 $B \approx 50\UuT$。B 正确；

CD. 由选项 A 可知测量地在北半球，而北半球地磁场指向北方斜向下，则第 2 次测量，测量 $B_{y} < 0$，故 $y$ 轴指向南方，第 3 次测量 $B_{x} > 0$ 故 $x$ 轴指向北方而 $y$ 轴则指向西方，C 正确、D 错误。

故选 BC。
}

\item
%% number: 6
%% paperName: 2022年全国高考乙卷第 6 题 6 分
%% typeId: 2
%% type: 多选题
%% chapter: 电场
%% point: 电场叠加
%% method: 无
%% score: 6
%% degree: 450
%% duplicateId: 0
%% body:
如图，两对等量异号点电荷 $+q$、$-q$（$q > 0$）固定于正方形的 4 个顶点上。L、N 是该正方形两条对角线与其内切圆的交点，O 为内切圆的圆心，M 为切点。则\xzanswer{AB}

\onepicture[w=4.2cm]{figs/06.png}

\fourchoices[answer=AB]
{L 和 N 两点处的电场方向相互垂直}
{M 点的电场方向平行于该点处的切线，方向向左}
{将一带正电的点电荷从 M 点移动到 O 点，电场力做正功}
{将一带正电的点电荷从 L 点移动到 N 点，电场力做功为零}

%% answer: AB
%% memo:
\memoanswer{
A. 两个正电荷在 N 点产生的场强方向由 N 指向 O，N 点处于两负电荷连线的中垂线上，则两负电荷在 N 点产生的场强方向由 N 指向 O，则 N 点的合场强方向由 N 指向 O，

同理可知，两个负电荷在 L 处产生的场强方向由 O 指向 L，L 点处于两正电荷连线的中垂线上，两正电荷在 L 处产生的场强方向由 O 指向 L，则 L 处的合场方向由 O 指向 L，由于正方形两对角线垂直平分，则 L 和 N 两点处的电场方向相互垂直，故 A 正确；

B. 正方形底边的一对等量异号电荷在 M 点产生的场强方向向左，而正方形上方的一对等量异号电荷在 M 点产生的场强方向向右，由于 M 点离上方一对等量异号电荷距离较远，则 M 点的场方向向左，故 B 正确；

C. 由图可知，M 和 O 点位于两等量异号电荷的等势线上，即 M 和 O 点电势相等，所以将一带正电的点电荷从 M 点移动到 O 点，电场力做功为零，故 C 错误；

D. 由图可知，L 点的电势低于 N 点电势，则将一带正电的点电荷从 L 点移动到 N 点，电场力做功不为零，故 D 错误。

故选 AB。
}

\item
%% number: 7
%% paperName: 2022年全国高考乙卷第 7 题 6 分
%% typeId: 2
%% type: 多选题
%% chapter: 运动和力
%% point: 牛顿定律中的图象
%% method: 无
%% score: 6
%% degree: 450
%% duplicateId: 0
%% body:
质量为 $1\Ukg$ 的物块在水平力 $F$ 的作用下由静止开始在水平地面上做直线运动，$F$ 与时间 $t$ 的关系如图所示。已知物块与地面间的动摩擦因数为 0.2，重力加速度大小取 $g = 10\Umsq$。则\xzanswer{AD}

\onepicture[w=7.0cm]{figs/07.png}

\fourchoices[answer=AD]
{$4\Us$ 时物块的动能为零}
{$6\Us$ 时物块回到初始位置}
{$3\Us$ 时物块的动量为 $12\Ukgms$}
{$0\sim 6\Us$ 时间内 $F$ 对物块所做的功为 $40\UJ$}

%% answer: AD
%% memo:
\memoanswer{
物块与地面间的摩擦力为 $f = \mu mg = 2\UN$

AC. 对物块从 $0\sim 3\Us$ 内由动量定理可知 $(F - f)t_{1} = mv_{3}$

即 $(4 - 2)\times 3 = 1\times v_{3}$

得 $v_{3} = 6\Ums$

$3\Us$ 时物块的动量为 $p = mv_{3} = 6\Ukgms$

设 $3\Us$ 后经过时间 $t$ 物块的速度减为 0，由动量定理可得 $-(F + f)t = 0 - mv_{3}$

即 $-(4 + 2)t = 0 - 1\times 6$

解得 $t = 1\Us$

所以物块在 $4\Us$ 时速度减为 0，则此时物块的动能也为 0，故 A 正确，C 错误；

B. $0\sim 3\Us$ 物块发生的位移为 $x_{1}$，由动能定理可得 $(F - f)x_{1} = \frac{1}{2}mv_{3}^{2}$

即 $(4 - 2)x_{1} = \frac{1}{2}\times 1\times 6^{2}$

得 $x_{1} = 9\Um$

$3\Us\sim 4\Us$ 过程中，对物块由动能定理可得 $-(F + f)x_{2} = 0 - \frac{1}{2}mv_{3}^{2}$

即 $-(4 + 2)x_{2} = 0 - \frac{1}{2}\times 1\times 6^{2}$

得 $x_{2} = 3\Um$

$4\Us\sim 6\Us$ 物块开始反向运动，物块的加速度大小为 $a = \frac{F - f}{m} = 2\Umsq$

发生的位移为 $x_{3} = \frac{1}{2}\times 2\times 2^{2}\Um = 4\Um < x_{1} + x_{2}$

即 $6\Us$ 时物块没有回到初始位置，故 B 错误；

D. 物块在 $6\Us$ 时的速度大小为 $v_{6} = 2\times 2\Ums = 4\Ums$

$0\sim 6\Us$ 拉力所做的功为 $W = (4\times 9 - 4\times 3 + 4\times 4)\UJ = 40\UJ$

故 D 正确。

故选 AD。
}

\item
%% number: 8
%% paperName: 2022年全国高考乙卷第 8 题 6 分
%% typeId: 2
%% type: 多选题
%% chapter: 电场
%% point: 带电粒子的圆周运动
%% method: 无
%% score: 6
%% degree: 450
%% duplicateId: 0
%% body:
一种可用于卫星上的带电粒子探测装置，由两个同轴的半圆柱形带电导体极板（半径分别为 $R$ 和 $R + d$）和探测器组成，其横截面如图\ref{2022全国乙卷08a}所示，点 O 为圆心。在截面内，极板间各点的电场强度大小与其到 O 点的距离成反比，方向指向 O 点。4 个带正电的同种粒子从极板间通过，到达探测器。不计重力。粒子 1、2 做圆周运动，圆的圆心为 O、半径分别为 $r_{1}$、$r_{2}$（$R < r_{1} < r_{2} < R + d$）；粒子 3 从距 O 点 $r_{2}$ 的位置入射并从距 O 点 $r_{1}$ 的位置出射；粒子 4 从距 O 点 $r_{1}$ 的位置入射并从距 O 点 $r_{2}$ 的位置出射，轨迹如图\ref{2022全国乙卷08b}中虚线所示。则\xzanswer{BD}

\twopicture[numA=a, numB=b, labelA=2022全国乙卷08a, labelB=2022全国乙卷08b, wA=5.0cm, wB=13.5cm]
{figs/08a.png}
{figs/08b.png}

\fourchoices[answer=BD]
{粒子 3 入射时的动能比它出射时的大}
{粒子 4 入射时的动能比它出射时的大}
{粒子 1 入射时的动能小于粒子 2 入射时的动能}
{粒子 1 入射时的动能大于粒子 3 入射时的动能}

%% answer: BD
%% memo:
\memoanswer{
C. 在截面内，极板间各点的电场强度大小与其到 O 点的距离成反比，可设为 $Er = k$

带正电的同种粒子 1、2 在均匀辐向电场中做匀速圆周运动，则有 $qE_{1} = m\frac{v_{1}^{2}}{r_{1}}$，$qE_{2} = m\frac{v_{2}^{2}}{r_{2}}$

可得 $\frac{1}{2}mv_{1}^{2} = \frac{qE_{1}r_{1}}{2} = \frac{qE_{2}r_{2}}{2}$

即粒子 1 入射时的动能等于粒子 2 入射时的动能，故 C 错误；

A. 粒子 3 从距 O 点 $r_{2}$ 的位置入射并从距 O 点 $r_{1}$ 的位置出射，做向心运动，电场力做正功，则动能增大，粒子 3 入射时的动能比它出射时的小，故 A 错误；

B. 粒子 4 从距 O 点 $r_{1}$ 的位置入射并从距 O 点 $r_{2}$ 的位置出射，做离心运动，电场力做负功，则动能减小，粒子 4 入射时的动能比它出射时的大，故 B 正确；

D. 粒子 3 做向心运动，有 $qE_{2} > m\frac{v_{3}^{2}}{r_{2}}$

可得 $\frac{1}{2}mv_{3}^{2} < \frac{qE_{2}r_{2}}{2} = \frac{1}{2}mv_{1}^{2}$

粒子 1 入射时的动能大于粒子 3 入射时的动能，故 D 正确。

故选 BD。
}

\item
%% number: 9
%% paperName: 2022年全国高考乙卷第 9 题 5 分
%% typeId: 3
%% type: 填空题
%% chapter: 直线运动
%% point: 打点计时器公式
%% method: 无
%% score: 5
%% degree: 650
%% duplicateId: 0
%% body:
用雷达探测一高速飞行器的位置。从某时刻（$t = 0$）开始的一段时间内，该飞行器可视为沿直线运动，每隔 $1\Us$ 测量一次其位置，坐标为 $x$，结果如下表所示：

\begin{table}[htbp]
\centering
\begin{tblr}{|c|c|c|c|c|c|c|c|}
\hline
$t/\Us[a]$ & 0 & 1 & 2 & 3 & 4 & 5 & 6 \\
\hline
$x/\Um[a]$ & 0 & 507 & 1094 & 1759 & 2505 & 3329 & 4233 \\
\hline
\end{tblr}
\end{table}

回答下列问题：
\begin{enumerate}
	\item
	根据表中数据可判断该飞行器在这段时间内近似做匀加速运动，判断的理由是：

	\tkanswer[5cm]{相邻 $1\Us$ 内的位移之差接近 $\Delta x = 80\Um$}；
	\item
	当 $x = 507\Um$ 时，该飞行器速度的大小 $v = $\tkanswer[1.6cm]{$547$} $\Ums$；
	\item
	这段时间内该飞行器加速度的大小 $a = $\tkanswer[1.6cm]{$79$} $\Umsq$（保留 2 位有效数字）。
\end{enumerate}

%% answer:
\jdanswer{
\begin{enumerate}
	\item
	相邻 $1\Us$ 内的位移之差接近 $\Delta x = 80\Um$

	\item
	$547$

	\item
	$79$

\end{enumerate}
}
%% memo:
\memoanswer{
（1）第 $1\Us$ 内的位移 $507\Um$，第 $2\Us$ 内的位移 $587\Um$，第 $3\Us$ 内的位移 $665\Um$，第 $4\Us$ 内的位移 $746\Um$，第 $5\Us$ 内的位移 $824\Um$，第 $6\Us$ 内的位移 $904\Um$，则相邻 $1\Us$ 内的位移之差接近 $\Delta x = 80\Um$，可知判断飞行器在这段时间内做匀加速运动；

（2）当 $x = 507\Um$ 时飞行器的速度等于 $0\sim 2\Us$ 内的平均速度，则 $v_{1} = \frac{1094}{2}\Ums = 547\Ums$

（3）根据 $a = \frac{x_{36} - x_{03}}{9T^{2}} = \frac{4233 - 2\times1759}{9\times1^{2}}\Umsq \approx 79\Umsq$
}

\item
%% number: 10
%% paperName: 2022年全国高考乙卷第 10 题 10 分
%% typeId: 4
%% type: 实验题
%% chapter: 电路
%% point: 伏安法测电阻
%% method: 无
%% score: 10
%% degree: 450
%% duplicateId: 0
%% body:
一同学探究阻值约为 $550\UO$ 的待测电阻 $R_{x}$ 在 $0\sim 5\UmA$ 范围内的伏安特性。可用器材有：电压表 V（量程为 $3\UV$，内阻很大），电流表 A（量程为 $1\UmA$，内阻为 $300\UO$），电源 $E$（电动势约为 $4\UV$，内阻不计），滑动变阻器 $R$（最大阻值可选 $10\UO$ 或 $1.5\UkO$），定值电阻 $R_{0}$（阻值可选 $75\UO$ 或 $150\UO$），开关 S，导线若干。
\begin{enumerate}
	\item
	要求通过 $R_{x}$ 的电流可在 $0\sim 5\UmA$ 范围内连续可调，在答题卡上将图\ref{2022全国乙卷10a}所示的器材符号连线，画出实验电路的原理图\tkanswer[2cm]{}；
	\item
	实验时，图\ref{2022全国乙卷10a}中的 $R$ 应选最大阻值为\tkanswer[1.8cm]{$10$} $\UO$（填“$10\UO$”或“$1.5\UkO$”）的滑动变阻器，$R_{0}$ 应选阻值为\tkanswer[1.8cm]{$75$} $\UO$（填“$75\UO$”或“$150\UO$”）的定值电阻；
	\item
	测量多组数据可得 $R_{x}$ 的伏安特性曲线。若在某次测量中，电压表、电流表的示数分别如图\ref{2022全国乙卷10b}和图\ref{2022全国乙卷10c}所示，则此时 $R_{x}$ 两端的电压为\tkanswer[1.6cm]{$2.30$} $\UV$，流过 $R_{x}$ 的电流为\tkanswer[1.6cm]{$4.20$} $\UmA$，此组数据得到的 $R_{x}$ 的阻值为\tkanswer[1.8cm]{$548$} $\UO$（保留 3 位有效数字）。
\end{enumerate}

\threepicture[numA=a, numB=b, numC=c, labelA=2022全国乙卷10a, labelB=2022全国乙卷10b, labelC=2022全国乙卷10c, wA=4.0cm, wB=6.0cm, wC=6.0cm]
{figs/10a.png}
{figs/10b.png}
{figs/10c.png}

%% answer:
\jdanswer{
\begin{enumerate}
	\item
	电路图见详解

	\item
	$10\UO$，$75\UO$

	\item
	$2.30$，$4.20$，$548$

\end{enumerate}
}
%% memo:
\memoanswer{
（1）电流表内阻已知，电流表与 $R_{0}$ 并联扩大电流表量程，进而准确测量通过 $R_{x}$ 的电流，电压表单独测量 $R_{x}$ 的电压；滑动变阻器采用分压式接法，电表从 0 开始测量，满足题中通过 $R_{x}$ 的电流从 $0\sim 5\UmA$ 连续可调，电路图如下

\onepicture[w=5.0cm]{figs/10_an.png}

（2）电路中 $R$ 应选最大阻值为 $10\UO$ 的滑动变阻器，方便电路的调节，测量效率高、实验误差小；

通过 $R_{x}$ 的电流最大为 $5\UmA$，需要将电流表量程扩大为原来的 5 倍，根据并联分流的规律示意图如下

\onepicture[w=6.5cm]{figs/10_an2.png}

根据并联分流，即并联电路中电流之比等于电阻的反比，可知 $\frac{4\UmA}{1\UmA}=\frac{300\UO}{R_{0}}$

解得 $R_{0}=75\UO$

（3）电压表每小格表示 $0.1\UV$，向后估读一位，即 $U = 3.20\UV$；电流表每小格表示 $0.02\UmA$，本位估读，即 $0.84\UmA$，电流表量程扩大 5 倍，所以通过 $R_{x}$ 的电流为 $I = 4.20\UmA$；

根据欧姆定律可知 $R_{x}=\frac{U}{I}=\frac{2.30}{4.20\times10^{-3}}\UO\approx548\UO$
}

\item
%% number: 11
%% paperName: 2022年全国高考乙卷第 11 题 12 分
%% typeId: 6
%% type: 计算题
%% chapter: 电磁感应
%% point: 法拉第电磁感应定律
%% method: 无
%% score: 12
%% degree: 650
%% duplicateId: 0
%% body:
如图，一不可伸长的细绳的上端固定，下端系在边长为 $l = 0.40\Um$ 的正方形金属框的一个顶点上。金属框的一条对角线水平，其下方有方向垂直于金属框所在平面的匀强磁场。已知构成金属框的导线单位长度的阻值为 $\lambda = 5.0\times10^{-3}\UOpm$；在 $t = 0$ 到 $t = 3.0\Us$ 时间内，磁感应强度大小随时间 $t$ 的变化关系为 $B(t) = 0.3 - 0.1t$（SI）。求：
\begin{enumerate}
	\item
	$t = 2.0\Us$ 时金属框所受安培力的大小；
	\item
	在 $t = 0$ 到 $t = 2.0\Us$ 时间内金属框产生的焦耳热。
\end{enumerate}

\onepicture[w=5.0cm, align=right]{figs/11.png}

%% answer:
\jdanswer{
\begin{enumerate}
	\item
	$0.04\sqrt{2}\UN$

	\item
	$0.016\UJ$

\end{enumerate}
}
%% memo:
\memoanswer{
（1）金属框的总电阻为 $R=4l\lambda=4\times0.4\times5\times10^{-3}\UO=0.008\UO$

金属框中产生的感应电动势为 $E=\frac{\Delta\Phi}{\Delta t}=\frac{\Delta B\times\frac{l^{2}}{2}}{\Delta t}=0.1\times\frac{1}{2}\times0.4^{2}\UV=0.008\UV$

金属框中的电流为 $I=\frac{E}{R}=1\UA$

$t=2.0\Us$ 时磁感应强度为 $B_{2}=(0.3-0.1\times2)\UT=0.1\UT$

金属框处于磁场中的有效长度为 $L=\sqrt{2}l$

此时金属框所受安培力大小为 $F_{A}=B_{2}IL=0.1\times1\times\sqrt{2}\times0.4\UN=0.04\sqrt{2}\UN$

（2）$0\sim 2.0\Us$ 内金属框产生的焦耳热为 $Q=I^{2}Rt=1^{2}\times0.008\times2\UJ=0.016\UJ$
}

\item
%% number: 12
%% paperName: 2022年全国高考乙卷第 12 题 20 分
%% typeId: 6
%% type: 计算题
%% chapter: 运动和力
%% point: 动量和能量
%% method: 无
%% score: 20
%% degree: 350
%% duplicateId: 0
%% body:
如图\ref{2022全国乙卷12a}，一质量为 $m$ 的物块 A 与轻质弹簧连接，静止在光滑水平面上：物块 B 向 A 运动，$t =0$ 时与弹簧接触，到 $t = 2t_{0}$ 时与弹簧分离，第一次碰撞结束，A、B 的 $v-t$ 图像如图\ref{2022全国乙卷12b}所示。已知从 $t = 0$ 到 $t = t_{0}$ 时间内，物块 A 运动的距离为 $0.36v_{0}t_{0}$。A、B 分离后，A 滑上粗糙斜面，然后滑下，与一直在水平面上运动的 B 再次碰撞，之后 A 再次滑上斜面，达到的最高点与前一次相同。斜面倾角为 $\theta$（$\sin\theta = 0.6$），与水平面光滑连接。碰撞过程中弹簧始终处于弹性限度内。求
\begin{enumerate}
	\item
	第一次碰撞过程中，弹簧弹性势能的最大值；
	\item
	第一次碰撞过程中，弹簧压缩量的最大值；
	\item
	物块 A 与斜面间的动摩擦因数。
\end{enumerate}

\twopicture[numA=a, numB=b, labelA=2022全国乙卷12a, labelB=2022全国乙卷12b, wA=6.5cm, wB=5.0cm, align=right]
{figs/12a.png}
{figs/12b.png}

%% answer:
\jdanswer{
\begin{enumerate}
	\item
	$0.6mv_{0}^{2}$

	\item
	$0.768v_{0}t_{0}$

	\item
	$0.45$

\end{enumerate}
}
%% memo:
\memoanswer{
（1）当弹簧被压缩最短时，弹簧弹性势能最大，此时 A、B 速度相等，即 $t = t_{0}$ 时刻，根据动量守恒定律 $m_{B}\cdot 1.2v_{0} = (m_{B} + m)v_{0}$

根据能量守恒定律 $E_{p\max} = \frac{1}{2}m_{B}(1.2v_{0})^{2} - \frac{1}{2}(m_{B} + m)v_{0}^{2}$

联立解得 $m_{B} = 5m$，$E_{p\max} = 0.6mv_{0}^{2}$

（2）解法一：同一时刻弹簧对 A、B 的弹力大小相等，根据牛顿第二定律 $F = ma$

可知同一时刻 $a_{A}=5a_{B}$

则同一时刻 A、B 的瞬时速度分别为 $v_{A}=a_{A}t$，$v_{B}=1.2v_{0}-\frac{a_{A}t}{5}$

根据位移等速度在时间上的累积可得 $s_{A}=v_{A}t$（累积），$s_{B}=v_{B}t$（累积）

又 $s_{A}=0.36v_{0}t_{0}$

解得 $s_{B}=1.128v_{0}t_{0}$

第一次碰撞过程中，弹簧压缩量的最大值 $\Delta s = s_{B} - s_{A} = 0.768v_{0}t_{0}$

解法二：B 接触弹簧后，压缩弹簧的过程中，A、B 动量守恒，有 $m_{B}\times 1.2v_{0} = 6mv_{0} = m_{B}v_{B} + mv_{A}$

对方程两边同时乘以时间 $\Delta t$，有 $6mv_{0}\Delta t = 5mv_{B}\Delta t + mv_{A}\Delta t$

$0-t_{0}$ 之间，根据位移等速度在时间上的累积，可得 $6mv_{0}t_{0}=5ms_{B}+ms_{A}$

将 $s_{A}=0.36v_{0}t_{0}$ 代入可得 $s_{B}=1.128v_{0}t_{0}$

则第一次碰撞过程中，弹簧压缩量的最大值 $\Delta s = s_{B} - s_{A} = 0.768v_{0}t_{0}$

（3）物块 A 第二次到达斜面的最高点与第一次相同，说明物块 A 第二次与 B 分离后速度大小仍为 $2v_{0}$，方向水平向右，设物块 A 第一次滑下斜面的速度大小为 $v_{A}$，设向左为正方向，根据动量守恒定律可得 $mv_{A}' - 5m\cdot 0.8v_{0} = m\cdot(-2v_{0}) + 5mv_{B}'$

根据能量守恒定律可得 $\frac{1}{2}mv_{A}^{2} + \frac{1}{2}\cdot 5m\cdot(0.8v_{0})^{2} = \frac{1}{2}m\cdot(-2v_{0})^{2} + \frac{1}{2}\cdot 5mv_{B}^{2}$

联立解得 $v_{A}' = v_{0}$

方法一：设在斜面上滑行的长度为 $L$，上滑过程，根据动能定理可得 $-mgL\sin\theta - \mu mgL\cos\theta = 0 - \frac{1}{2}m(2v_{0})^{2}$

下滑过程，根据动能定理可得 $mgL\sin\theta - \mu mgL\cos\theta = \frac{1}{2}mv_{0}^{2} - 0$

联立解得 $\mu = 0.45$

方法二：根据牛顿第二定律，可以分别计算出滑块 A 上滑和下滑时的加速度，$mg\sin\theta + \mu mg\cos\theta = ma_{\text{上}}$，$mg\sin\theta - \mu mg\cos\theta = ma_{\text{下}}$

上滑时末速度为 0，下滑时初速度为 0，由匀变速直线运动的位移速度关系可得 $2a_{\text{上}}x = (2v_{0})^{2} - 0$，$2a_{\text{下}}x = v_{A}'^{2} = v_{0}^{2}$

联立可解得 $\mu = 0.45$
}

\item
%% number: 13
%% paperName: 2022年全国高考乙卷第 13 题 10 分
%% typeId: 6
%% type: 计算题
%% chapter: 气体性质
%% point: 盖吕萨克定律
%% method: 无
%% score: 10
%% degree: 450
%% duplicateId: 0
%% body:
（1）一定量的理想气体从状态 a 经状态 b 变化到状态 c，其过程如 $T-V$ 图上的两条线段所示，则气体在\xzanswer{ABD}

\onepicture[w=4.5cm]{figs/13a.png}

\fivechoices[answer=ABD]
{状态 a 处的压强大于状态 c 处的压强}
{由 a 变化到 b 的过程中，气体对外做功}
{由 b 变化到 c 的过程中，气体的压强不变}
{由 a 变化到 b 的过程中，气体从外界吸热}
{由 a 变化到 b 的过程中，从外界吸收的热量等于其增加的内能}

（2）如图，一竖直放置的汽缸由两个粗细不同的圆柱形筒组成，汽缸中活塞Ⅰ和活塞Ⅱ之间封闭有一定量的理想气体，两活塞用一轻质弹簧连接，汽缸连接处有小卡销，活塞Ⅱ不能通过连接处。活塞Ⅰ、Ⅱ的质量分别为 $2m$、$m$，面积分别为 $2S$、$S$，弹簧原长为 $l$。初始时系统处于平衡状态，此时弹簧的伸长量为 $0.1l$，活塞Ⅰ、Ⅱ到汽缸连接处的距离相等，两活塞间气体的温度为 $T_{0}$。已知活塞外大气压强为 $p_{0}$，忽略活塞与缸壁间的摩擦，汽缸无漏气，不计弹簧的体积。
\begin{enumerate}
	\item
	求弹簧的劲度系数；
	\item
	缓慢加热两活塞间的气体，求当活塞Ⅱ刚运动到汽缸连接处时，活塞间气体的压强和温度。
\end{enumerate}

\onepicture[w=2.6cm, align=right]{figs/13b.png}

%% answer:
\jdanswer{
\begin{enumerate}
	\item
	$k = \frac{40mg}{l}$

	\item
	$p_{2} = p_{0} + \frac{3mg}{S}$，$T_{2} = \frac{4}{3}T_{0}$

\end{enumerate}
}
%% memo:
\memoanswer{
AC. 根据理想气体状态方程可知 $T=\frac{p}{nR}\cdot V$

即 $T-V$ 图像的斜率为 $\frac{p}{nR}$，故有 $p_{a}=p_{b}>p_{c}$

故 A 正确，C 错误；

B. 理想气体由 a 变化到 b 的过程中，因体积增大，则气体对外做功，故 B 正确；

DE. 理想气体由 a 变化到 b 的过程中，温度升高，则内能增大，由热力学第一定律有 $\Delta U = Q + W$

而 $\Delta U > 0$，$W < 0$，则有 $\Delta U = Q - |W|$

可得 $Q > 0$，$Q > \Delta U$

即气体从外界吸热，且从外界吸收的热量大于其增加的内能，故 D 正确，E 错误；

故选 ABD。

（1）设封闭气体的压强为 $p_{1}$，对两活塞和弹簧的整体受力分析，由平衡条件有 $mg + p_{0}\cdot 2S + 2mg + p_{1}S = p_{0}S + p_{1}\cdot 2S$

解得 $p_{1}=p_{0}+\frac{3mg}{S}$

对活塞Ⅰ由平衡条件有 $2mg + p_{0}\cdot 2S + k\cdot 0.1l = p_{1}\cdot 2S$

解得弹簧的劲度系数为 $k=\frac{40mg}{l}$

（2）缓慢加热两活塞间的气体使得活塞Ⅱ刚运动到汽缸连接处时，对两活塞和弹簧的整体由平衡条件可知，气体的压强不变依然为 $p_{2}=p_{1}=p_{0}+\frac{3mg}{S}$

即封闭气体发生等压过程，初末状态的体积分别为 $V_{1}=\frac{1.1l}{2}\times2S+\frac{1.1l}{2}\times S=\frac{3.3lS}{2}$，$V_{2}=l_{2}\cdot2S$

由气体的压强不变，则弹簧的弹力也不变，故有 $l_{2}=1.1l$

有等压方程可知 $\frac{V_{1}}{T_{0}}=\frac{V_{2}}{T_{2}}$

解得 $T_{2}=\frac{4}{3}T_{0}$
}

\item
%% number: 14
%% paperName: 2022年全国高考乙卷第 14 题 10 分
%% typeId: 6
%% type: 计算题
%% chapter: 光学
%% point: 全反射
%% method: 无
%% score: 10
%% degree: 550
%% duplicateId: 0
%% body:
（1）介质中平衡位置在同一水平面上的两个点波源 $S_{1}$ 和 $S_{2}$，二者做简谐运动的振幅相等，周期均为 $0.8\Us$。当 $S_{1}$ 过平衡位置向上运动时，$S_{2}$ 也过平衡位置向上运动。若波速为 $5\Ums$，则由 $S_{1}$ 和 $S_{2}$ 发出的简谐横波的波长均为\tkanswer[1.6cm]{$4$} $\Um$。P 为波源平衡位置所在水平面上的一点，与 $S_{1}$、$S_{2}$ 平衡位置的距离均为 $10\Um$，则两波在 P 点引起的振动总是相互\tkanswer[1.6cm]{加强}的；当 $S_{1}$ 恰好在平衡位置向上运动时，平衡位置在 P 处的质点\tkanswer[1.6cm]{向下}运动。

（2）一细束单色光在三棱镜 $ABC$ 的侧面 $AC$ 上以大角度由 $D$ 点入射（入射面在棱镜的横截面内），入射角为 $i$，经折射后射至 $AB$ 边的 $E$ 点，如图所示，逐渐减小 $i$，$E$ 点向 $B$ 点移动，当 $\sin i = \frac{1}{6}$ 时，恰好没有光线从 $AB$ 边射出棱镜，且 $DE = DA$。求棱镜的折射率。

\onepicture[w=5.0cm, align=right]{figs/14.png}

%% answer:
\jdanswer{$n = 1.5$}
%% memo:
\memoanswer{
（1）因周期 $T = 0.8\Us$，波速为 $v = 5\Ums$，则波长为 $\lambda = vT = 4\Um$。

（2）因两波源到 P 点的距离之差为零，且两振源振动方向相同，则 P 点的振动是加强的；

（3）因 $S_{1}P = 10\Um = 2.5\lambda$，则当 $S_{1}$ 恰好的平衡位置向上运动时，平衡位置在 P 点的质点向下振动。

因为当 $\sin i = \frac{1}{6}$ 时，恰好没有光线从 AB 边射出，可知光线在 E 点发生全反射，设临界角为 $C$，则 $\sin C = \frac{1}{n}$

由几何关系可知，光线在 D 点的折射角为 $r = 90^\circ - 2C$

则 $\frac{\sin i}{\sin r} = n$

联立可得 $n = 1.5$

\onepicture[w=5.5cm]{figs/14_an.png}
}

\end{enumerate}

\end{document}
